% Generated draft: requires mathematical and source review.
\documentclass[11pt]{article}
\usepackage[margin=25mm]{geometry}
\usepackage{fontspec}
\setmainfont{Latin Modern Roman}
\usepackage{amsmath,amssymb,mathtools,unicode-math}
\setmathfont{Latin Modern Math}
\usepackage{enumitem,needspace,xurl,hyperref}
\setlength{\parindent}{0pt}
\setlength{\parskip}{5pt}
\setlength{\emergencystretch}{3em}
\newcommand{\N}{\mathbb{N}}
\newcommand{\Z}{\mathbb{Z}}
\newcommand{\Q}{\mathbb{Q}}
\newcommand{\R}{\mathbb{R}}
\newcommand{\C}{\mathbb{C}}
\newcommand{\F}{\mathbb{F}}
\newcommand{\A}{\mathbb{A}}
\newcommand{\PP}{\mathbb{P}}
\newcommand{\E}{\mathbb{E}}
\newcommand{\eps}{\varepsilon}
\newcommand{\dd}{\,\mathrm{d}}
\begin{document}
\section*{UnsolvedMath ID: 1200043}
\textbf{Source problem code:} catalogue entry
\subsection{Definitions and mathematical statement}
In the usual cost framework, let $\Gamma$ be a countable discrete group acting by probability-measure-preserving transformations on a standard probability space $(X,\mu)$. Its orbit equivalence relation is
\[
R_{\Gamma}=\{(x,\gamma x):x\in X,\ \gamma\in\Gamma\},
\]
understood modulo null sets. A graphing of $R_{\Gamma}$ is a countable family of measure-preserving measurable partial bijections $\phi_j:A_j\to B_j$ whose graphs generate that equivalence relation modulo null sets. Define
\[
\operatorname{cost}(\Gamma,X)=\inf_{(\phi_j)}\sum_j\mu(A_j).
\]
For $H\leq\Gamma$, $\operatorname{cost}(H,X)$ uses the restricted action on the same probability space, not a new action or a changed measure.

\textbf{Question.} For every such action and every subgroup $H$ of finite index $d=[\Gamma:H]$, does
\[
\operatorname{cost}(H,X)-1
=d\big(\operatorname{cost}(\Gamma,X)-1\big)
\]
hold? Costs may be infinite, with the usual extended-real convention for this formula. The source says only probability-measure-preserving and does not explicitly require an essentially free action, meaning $\mu(\{x:\gamma x=x\})=0$ for every $\gamma\ne1$; neither ergodicity nor atomlessness is stated. These hypotheses cannot be silently inserted.

The supplied research summary reports the scaled identity in fixed-price settings, citing Gaboriau. A group has fixed price when the costs of all its essentially free probability-measure-preserving actions coincide. Since that condition concerns essentially free actions, the reported progress does not by itself validate the literal formula for arbitrary nonfree actions. The summary links the unrestricted intended problem to the open fixed-price problem rather than claiming an unconditional theorem.
\subsection{Short English statement}
Does restricting a probability-preserving action to a finite-index subgroup multiply cost minus one by the index? The source’s missing freeness hypotheses and its fixed-price qualification must be kept distinct.
\subsection{Sources}
\begin{itemize}
\item M. Ab\'ert, \emph{Some questions} (2010), Question 43, \url{https://www.renyi.hu/~abert/questions.pdf}; manual transcription, accessed 2026-07-29, with current status not independently verified.
\item Gaboriau (2000), cited by the research summary for fixed-price progress. No publication title, theorem number, or further bibliographic identifier is supplied.
\item Local registry \texttt{data/unsolvedmath/problems.json}, permanent ID $1200043$, source code catalogue entry, \emph{Some Questions --- Question 43}. The record labels this partially solved and PARTIAL-PROGRESS, with NEEDS\_REVIEW source-status and rights qualifications.
\end{itemize}
\end{document}
